%%%%%%%%%%%%%%%%%%%%%%%%%%%%%%%%%%%%%%%%%%%%%%%%%%%%%%%%%%%%%%%%%%%%%%
% Lớp 10 | Xác suất | Xác suất của biến cố
% Dạng 0D0?2-2: Liên quan xúc xắc, đồng tiền (PP liệt kê)
% Số câu: 43  (sinh tự động bởi tach_cau.py; từ nay thầy sửa trực tiếp file này)
%%%%%%%%%%%%%%%%%%%%%%%%%%%%%%%%%%%%%%%%%%%%%%%%%%%%%%%%%%%%%%%%%%%%%%

%%%%%%%%%%%%%%%% TRẮC NGHIỆM 4 PHƯƠNG ÁN %%%%%%%%%%%%%%%%

% nguồn: MaHoaNH_TracNghiem_Lop10
\begin{ex}%[0D0N2-2]%[ma:b37877b5]
	Gieo một con súc sắc. Xác suất để mặt chấm  chẵn xuất hiện là
	 \choice
		{$0{,}2$} 
		 {$0{,}3$}
		 {$0{,}4$}
		{\True $0{,}5$} 
	\loigiai{Không gian mẫu của phép thử trên  là $\Omega = \left\{1;2;3;4;5;6\right\}$ nên $n(\Omega)=6$.\\
		Gọi $A$ là biến cố: \lq \lq  Mặt chấm chẵn xuất hiện \rq \rq. Khi đó $A =\left\{2;4;6\right\}$ nên $n(A)=3$.\\
	Vậy $\mathrm{P}  \left(A\right)=\dfrac{n(A)}{n(\Omega)}=\dfrac{1}{2}=0{,}5$.}
\end{ex}

% nguồn: MaHoaNH_TracNghiem_Lop10
\begin{ex}%[0D0N2-2]%[ma:454da530]
	Tung một đồng xu cân đối và đồng chất $3$ lần liên tiếp. Tính xác  suất của biến cố $A:$ \lq \lq Trong $3$ lần tung có ít nhất $2$ lần xuất hiện mặt ngửa\rq \rq.
	 \choice
		 {$\dfrac{3}{8}$}
		{\True  $\dfrac{1}{2}$}
		{$\dfrac{5}{8}$} 
		{$\dfrac{2}{3}$} 
	\loigiai{Kí hiệu $S$ nếu tung được mặt  sấp, $N$ nếu tung được mặt ngửa.\\
		Ta có số phần tử của không gian mẫu  là $n(\Omega)=2^3=8$.\\
		Vì $A =\left\{NNN;NNS;NSN;SNN\right\}$  nên $n(A)=4$ và $\mathrm{P} \left(A\right)=\dfrac{4}{8}=\dfrac{1}{2}$.
	 }
\end{ex}

% nguồn: MaHoaNH_TracNghiem_Lop10
\begin{ex}%[0D0N2-2]%[ma:4805bc2a]
	Gieo đồng thời hai con xúc xắc cân đối và đồng chất. Tính xác suất của biến cố \lq \lq  Tổng số chấm trên hai mặt xuất hiện nhỏ hơn $5$\rq \rq.
	\choice 
		 {\True $\dfrac{1}{6}$}
		{$\dfrac{1}{3}$} 
		 {$\dfrac{5}{18}$}
		 {$\dfrac{5}{36}$}
	\loigiai{Ta có số  phần tử của không gian mẫu là $n \left(\Omega\right)=6\cdot 6 =36$.\\
		Gọi $A$ là biến cố: \lq \lq Tổng số chấm trên hai mặt xuất hiện nhỏ  hơn $5$\rq \rq.\\
		 Ta có $A=\left\{\left(1;1\right);\left(1;2\right);\left(1;3\right);\left(2;1\right);\left(2;2\right);\left(3;1\right)\right\}$.\\
		Xác  suất của biến cố $A$ là $\mathrm{P}(A)=\dfrac{6}{36}=\dfrac{1}{6}$.
	 }
\end{ex}

% nguồn: MaHoaNH_TracNghiem_Lop10
\begin{ex}%[0D0N2-2]%[ma:6c04364e]
	Gieo $1$ đồng xu cân đối và đồng chất $3$ lần. Xác suất để có đúng $2$ lần  gieo xuất hiện mặt $S$ là
	\choice 
		 {$\dfrac{3}{4}$}
		{\True  $\dfrac{3}{8}$}
		 {$\dfrac{1}{2}$}
		{$\dfrac{1}{4}$} 
	\loigiai{ 
		Không  gian mẫu $\Omega =\left\{SSS;NNN;SSN;SNS;NSS;NNS;NSN;SNN\right\}$.\\
		Số phần tử  của không gian mẫu $n(\Omega)=8$.\\
		Gọi biến cố $A$: \lq \lq Có đúng $2$ lần gieo xuất hiện mặt sấp\rq \rq \, nên ta có  $n(A)=3$.\\
		Xác suất của biến cố  $A$ là $\mathrm{P}(A)=\dfrac{3}{8}$.
	} 
\end{ex}

% nguồn: MaHoaNH_TracNghiem_Lop10
\begin{ex}%[0D0N2-2]%[ma:9d7c1197]
	Từ một hộp chứa $5$ quả cầu trắng và $3$ quả cầu đen, lấy ngẫu nhiên $2$ quả. Xác suất để lấy được cả  hai quả trắng là
	\choice 
		 {\True $\dfrac{5}{14}$}
		 {$\dfrac{5}{28}$}
		{$\dfrac{10}{56}$} 
		 {$\dfrac{20}{28}$}
	\loigiai{ 
		Tổng số quả là $5  + 3 = 8$.\\
		 Số cách chọn $2$ quả bất kỳ là $\mathrm{C}^2_8 = 28\Rightarrow$ số phần tử của không gian mẫu là $ n(\Omega)=28$.\\
		Gọi $A$ là biến cố \lq\lq Lấy được hai quả cầu  trắng\rq\rq.\\
		Số cách chọn $2$ quả trắng là $\mathrm{C}^2_5  = 10\Rightarrow n(A)=10$.\\
		Vậy  xác suất của biến cố $A$ là
		\[ 
		\mathrm{P}(A) = \dfrac{10}{28} = \dfrac{5}{14}. 
		 \]
	 }
\end{ex}

% nguồn: MaHoaNH_TracNghiem_Lop10
\begin{ex}%[0D0H2-2]%[ma:9c14ce3a]
	Gieo một con xúc xắc cân đối và đồng chất hai lần. Xác suất để  không một lần xuất hiện mặt sáu chấm là
	\choice 
		 {$\dfrac{12}{36}$}
		{\True  $\dfrac{25}{36}$}
		 {$\dfrac{6}{36}$}
		{$\dfrac{8}{36}$} 
	\loigiai{Số  phần tử của không gian mẫu là $n(\Omega)=6 \cdot 6=36$.\\
		Gọi $A$  là biến cố: \lq \lq Không có lần nào xuất hiện mặt sáu chấm\rq \rq.\\
	Ta có $n(A)=5 \cdot 5=25$ nên  $\mathrm{P}(A)=\dfrac{25}{36}$.}
\end{ex}

% nguồn: MaHoaNH_TracNghiem_Lop10
\begin{ex}%[0D0H2-2]%[ma:6cd6cc07]
	Gieo một con súc xắc cân đối, đồng chất một lần. Xác suất để xuất hiện mặt  chẵn?
	 \choice
		{\True  $\dfrac{1}{2}$}
		{$\dfrac{1}{4}$} 
		{$\dfrac{1}{6}$} 
		{$\dfrac{1}{3}$} 
	 \loigiai{
		Vì con súc xắc có $3$ mặt chẵn  trên tổng $6$ mặt nên xác suất xuất hiện mặt chẵn là $\dfrac{3}{6}=\dfrac{1}{2}$.
	 }
\end{ex}

% nguồn: MaHoaNH_TracNghiem_Lop10
\begin{ex}%[0D0H2-2]%[ma:0f3c9ed4]
	Gieo một đồng tiền liên tiếp $3$ lần. Xác suất của biến cố $A\colon$ \lq\lq Kết quả của  3 lần gieo là như nhau\rq\rq\, là
	 \choice
		 {$\mathrm{P}(A) = \dfrac{1}{2}$}
		{$\mathrm{P}(A)  = \dfrac{3}{8}$}
		 {$\mathrm{P}(A) = \dfrac{7}{8}$}
		 {\True $\mathrm{P}(A) = \dfrac{1}{4}$}
	 \loigiai{
		Vì $A\colon$ \lq\lq Kết quả của 3 lần gieo là như nhau\rq\rq\, nên  $A = \left\lbrace SSS; NNN\right\rbrace \Rightarrow n(A) = 2$ và $n(\Omega) = 2^3 = 8$. Xác suất của biến cố $A\colon$ là $\mathrm{P}(A) = \dfrac{n(A)}{n(\Omega)} = \dfrac{2}{8} = \dfrac{1}{4}$.
	} 
\end{ex}

% nguồn: MaHoaNH_TracNghiem_Lop10
\begin{ex}%[0D0H2-2]%[ma:5145e4c2]
	Gieo một con xúc xắc $3$ lần liên tiếp. Tính xác suất sao cho mặt năm chấm xuất hiện ở lần thứ hai. 
	 \choice
		{\True $\dfrac{1}{6}$} 
		 {$\dfrac{5}{36}$}
		{$\dfrac{1}{120}$} 
		{$\dfrac{5}{6}$} 
	\loigiai{ 
		Ta có số phần  tử của không gian mẫu là $n(\Omega)=6^3=216$.\\
		Gọi $A$ là biến cố: \lq \lq Mặt năm chấm xuất hiện  ở lần thứ hai \rq \rq.\\
		$A$ xảy  ra khi lần thứ hai xuất hiện mặt năm chấm, còn lần thứ nhất và lần thứ ba có thể xuất hiện các mặt bất kì trong $6$ mặt. Do đó $n(A)=6 \cdot 1 \cdot 6 =36$.\\
	Vậy xác suất của biến cố $A$  là $\mathrm{P}(A)=\dfrac{36}{216}=\dfrac{1}{6}$.}
\end{ex}

% nguồn: MaHoaNH_TracNghiem_Lop10
\begin{ex}%[0D0H2-2]%[ma:7e9d1363]
	Gieo ngẫu nhiên $2$ con xúc xắc cân đối đồng chất. Tìm xác suất của biến cố \lq \lq Hiệu số chấm xuất hiện  trên $2$ con xúc xắc bằng $2$\rq \rq.
	\choice 
		{$\dfrac{5}{18}$} 
		{$\dfrac{1}{9}$} 
		{\True $\dfrac{2}{9}$} 
		{$\dfrac{5}{6}$} 
	\loigiai{Số phần tử của  không gian mẫu là $n(\Omega)=6 \cdot 6=36$.\\
		Gọi $A$ là  biến cố thỏa mãn yêu cầu bài toán.\\
		 Khi đó $A=\left\{\left(1;3\right);\left(3;1\right);\left(2;4\right);\left(4;2\right);\left(3;5\right);\left(5;3\right);\left(4;6\right);\left(6;4\right)\right\}$ nên $n(A)=8$.\\
		Vậy $\mathrm{P}(A)=\dfrac{8}{36}=\dfrac{2}{9}$. 
	 }
\end{ex}

% nguồn: MaHoaNH_TracNghiem_Lop10
\begin{ex}%[0D0H2-2]%[ma:4e41bad6]
	Tung một đồng xu cân đối và đồng chất ba lần liên tiếp. Xác suất của  biến cố \lq \lq Mặt sấp xuất hiện ít nhất hai lần\rq \rq \ là
	 \choice
		 {$\dfrac{1}{4}$}
		{$\dfrac{1}{3}$} 
		 {\True $\dfrac{1}{2}$}
		 {$\dfrac{3}{4}$}
	\loigiai{ 
		Tung đồng xu  $3$ lần nên $n(\Omega)=8$.\\
		Gọi $A$ là biến cố \lq \lq Mặt sấp xuất hiện ít nhất hai lần \rq \rq. Khi  đó $n(A)=4$.\\
		Vậy $\mathrm{P}(A)=\dfrac{n(A)}{n(\Omega)=\dfrac{4}{8}}=\dfrac{1}{2}$. 
	} 
\end{ex}

% nguồn: MaHoaNH_TracNghiem_Lop10
\begin{ex}%[0D0H2-2]%[ma:ad260314]
	Gieo một con súc sắc cân đối và đồng chất,  với không gian mẫu $\Omega =\left\{1,2,3,4,5,6 \right\}$, gọi $E$ là biến cố \lq\lq mặt có số chấm chẵn xuất hiện\rq\rq. Kết quả nào sau đây là xác suất của $E$?
	 \choice
		 {$\dfrac{1}{3}$}
		 {\True $\dfrac{1}{2}$}
		 {$\dfrac{1}{6}$}
		{$\dfrac{3}{5}$} 
	\loigiai{ 
		Số phần tử của không gian mẫu là $n(\Omega)=6$.\\ 
		Biến cố \lq\lq mặt có số chấm chẵn xuất hiện\rq\rq~  là $E=\left\{2;4;6 \right\}$, $n(E)=3$.\\
		Suy  ra xác suất là $\mathrm{P}(E)=\dfrac{3}{6}=\dfrac{1}{2}$.
	} 
\end{ex}

% nguồn: MaHoaNH_TracNghiem_Lop10
\begin{ex}%[0D0H2-2]%[ma:dfc3626e]
	Gieo con súc sắc hai lần.  Biến cố $A$ là biến cố để sau hai lần gieo có ít nhất một mặt $6$ chấm
	\choice 
		{$A=\left\{\left(6;1\right),\left(6;2\right),\left(6;3\right),\left(6;4\right),\left(6;5\right)\right\}$} 
		 {$A=\left\{\left(1;6\right),\left(2;6\right),\left(3;6\right),\left(4;6\right),\left(5;6\right),\left(6;6\right)\right\}$}
		{\True $A=\left\{\left(1;6\right),\left(2;6\right),\left(3;6\right),\left(4;6\right),\left(5;6\right),\left(6;6\right),\left(6;1\right),\left(6;2\right),\left(6;3\right),\left(6;4\right),\left(6;5\right)\right\}$} 
		{$A=\left\{\left(1;6\right),\left(2;6\right),\left(3;6\right),\left(4;6\right),\left(5;6\right)\right\}$} 
	 \loigiai{
		Để sau 2 lần có ít nhất  1 mặt $6$ chấm ta có $2$ trường hợp sau:\\
		\textbf{Trường hợp 1:} Lần $1$ xuất  hiện mặt $6$ chấm, lần $2$ tùy ý.\\
		\textbf{Trường hợp 2:}  Lần $2$ xuất hiện mặt $6$ chấm, lần $1$ tùy ý.
	 }
\end{ex}

% nguồn: MaHoaNH_TracNghiem_Lop10
\begin{ex}%[0D0H2-2]%[ma:ae4cbe79]
	Tung một  đồng xu cân đối và đồng chất một lần. Xác suất của biến cố: \lq \lq Kết quả tung đồng xu là mặt ngửa\rq \rq \ là
	\choice 
		 {$\dfrac{1}{3}$}
		 {$\dfrac{3}{4}$}
		{$\dfrac{1}{4}$} 
		{\True  $\dfrac{1}{2}$}
	\loigiai{ 
		 Ta có $n(\Omega)=2$.\\
		Gọi $A$ là biến  cố: \lq \lq Tung đồng xu là mặt ngửa\rq \rq. Khi đó $n(A)=\dfrac{1}{2}$.\\
		Vậy $\mathrm{P}(A)=\dfrac{n(A)}{n(\Omega)}=\dfrac{1}{2}$. 
	 }
\end{ex}

% nguồn: MaHoaNH_TracNghiem_Lop10
\begin{ex}%[0D0H2-2]%[ma:303f3464]
	Gieo ngẫu nhiên một  con xúc xắc cân đối và đồng chất một lần. Tính xác suất để mặt $3$ chấm xuất hiện.
	\choice 
		 {$\dfrac{5}{6}$}
		{\True  $\dfrac{1}{6}$}
		{$\dfrac{1}{2}$} 
		 {$\dfrac{1}{3}$}
	\loigiai{ 
		Gieo ngẫu nhiên một con xúc xắc cân đối một lần, số phần tử của  không gian mẫu là $n(\Omega)=6$.\\
		Xét biến cố  $A$ \lq\lq Mặt $3$ chấm xuất hiện\rq\rq.\\
		Vậy  số khả năng thuận lợi cho $A$ là $n(A)=1$.\\
		Xác  suất của biến cố $A$ là $\mathrm{P}(A)=\dfrac{n(A)}{n(\Omega)}=\dfrac{1}{6}$.
	 }
\end{ex}

% nguồn: MaHoaNH_TracNghiem_Lop10
\begin{ex}%[0D0H2-2]%[ma:b976c0dc]
	Gieo ba con xúc  xắc cân đối và đồng chất. Tính xác suất để số chấm xuất hiện trên ba con xúc xắc đều bằng nhau.
	\choice 
		 {$\dfrac{1}{120}$}
		{$\dfrac{1}{720}$} 
		{\True  $\dfrac{1}{36}$}
		 {$\dfrac{1}{20}$}
	 \loigiai{
		Tập hợp  $\Omega$ có $216$ phần tử.\\
		Gọi $D$ là biến  cố \lq\lq số chấm xuất hiện trên ba con xúc xắc đều bằng nhau\rq\rq.\\
		Các  kết quả thuận lợi $(1;1;1)$, $(2;2;2)$, $(3;3;3)$, $(4;4;4)$, $(5;5;5)$, $(6;6;6)$.\\
		Tập hợp $D$ có $6$  phần tử.\\
		Vậy xác suất biến cố $D$ là $\mathrm{P}(D) = \dfrac{n(D)}{n(\Omega)} = \dfrac{6}{216} =  \dfrac{1}{36}$.
	 }
\end{ex}

% nguồn: MaHoaNH_TracNghiem_Lop10
\begin{ex}%[0D0H2-2]%[ma:0b1f1f9b]
	Gieo hai đồng xu $A$ và $B$ một cách độc lập. Đồng xu $A$ chế tạo cân đối. Đồng xu $B$ chế tạo không cân đối nên xác suất xuất hiện mặt sấp gấp $3$ lần xác suất xuất hiện mặt ngửa. Tính xác suất để khi gieo hai đồng xu cùng lúc được kết quả $1$ sấp và $1$ ngửa. 
	 \choice
		 {$\dfrac{1}{4}$}
		 {\True $\dfrac{1}{2}$}
		 {$\dfrac{3}{5}$}
		{$\dfrac{3}{4}$} 
		Lời 
	 \loigiai{
		Gieo đồng xu $B$, xác suất xuất hiện mặt ngửa là $\dfrac{1}{4}$ và xác suất xuất  hiện mặt sấp là $\dfrac{3}{4}$.\\
		Gieo đồng xu $A$, xác suất xuất hiện  mặt ngửa là $\dfrac{1}{2}$ và xác suất xuất hiện mặt sấp là $\dfrac{1}{2}$.\\
		Khi gieo hai đồng xu, để kết quả là $1$ sấp và $1$ ngửa thì có $2$ trường  hợp $A$ sấp $B$ ngửa hoặc $A$ ngửa $B$ sấp.\\
		Do đó  ta có xác suất $\mathrm{P}=\dfrac{1}{2} \cdot \dfrac{1}{4}+\dfrac{1}{2} \cdot \dfrac{3}{4}=\dfrac{1}{2}$.
	} 
\end{ex}

% nguồn: MaHoaNH_TracNghiem_Lop10
\begin{ex}%[0D0H2-2]%[ma:77517202]
	Xét phép thử $T$: \lq\lq Tung một đồng xu cân đối và đồng chất $2$ lần\rq\rq. Xác suất để lần thứ hai xuất hiện mặt $S$  là
	 \choice
		 {$0$}
		{$\dfrac{1}{4}$} 
		{\True $\dfrac{1}{2}$} 
		 {$1$}
	\loigiai{ 
		Không gian mẫu  là $\Omega=\{SS,SN,NS,NN\}$.\\
		Gọi $A$ là biến cố: \lq\lq Lần thứ hai xuất hiện mặt  sấp\rq\rq\,.\\
		 Ta có $A=\{SS,NS\}$. \\
		 $n(A)=2$.\\
		Xác suất  của biến cố $A$ là $\mathrm{P}(A)=\dfrac{n(A)}{n(\Omega)}=\dfrac{2}{4}=\dfrac{1}{2}$.
	} 
\end{ex}

% nguồn: MaHoaNH_TracNghiem_Lop10
\begin{ex}%[0D0H2-2]%[ma:01445f46]
	Gieo $2$ con súc sắc cân đối và đồng chất. Xác suất để tổng  số chấm xuất hiện trên hai mặt của $2$ con súc sắc đó không vượt quá $5$ là
	 \choice
		{$\dfrac{2}{3}$} 
		{$\dfrac{7}{18}$} 
		 {$\dfrac{8}{9}$}
		{\True $\dfrac{5}{18}$} 
	 \loigiai{
		Ta có $n(\Omega) = 36$.\\ 
		Gọi $A\colon$ \lq\lq tổng số  chấm xuất hiện trên hai mặt của $2$ con súc sắc đó không vượt quá $5$ \rq\rq.\\
		Khi đó $ A = \left\lbrace (1;1),(1;2),(2;1),(1;3),(3;1),(2;2),(1;4),(4;1),(2;3),(3;2)  \right\rbrace \Rightarrow n(A) = 10$.\\
		Xác suất để tổng số chấm xuất hiện trên hai mặt của $2$ con  súc sắc đó không vượt quá $5$ là
		 $$ \mathrm{P}(A) = \dfrac{n(A)}{n(\Omega)} = \dfrac{10}{36} = \dfrac{5}{18}. $$
	} 
\end{ex}

% nguồn: MaHoaNH_TracNghiem_Lop10
\begin{ex}%[0D0H2-2]%[ma:30fd5e6d]
	Gieo một đồng xu cân đối  đồng chất ba lần, xác suất để trong ba lần gieo có đúng hai lần xảy ra mặt sấp bằng
	\choice 
		 {$\dfrac{1}{8}$}
		{$\dfrac{3}{16}$} 
		{$\dfrac{3}{4}$} 
		{\True  $\dfrac{3}{8}$}
	 \loigiai{
		\begin{itemize}
			\item Số phần  tử của không gian mẫu là $n(\Omega)=2^3=8$.
			\item Gọi $A$ là biến cố: \lq\lq Ba lần gieo có đúng hai lần xảy ra  mặt sấp\rq\rq.\\
			Ta có $A=\{SSN;SNS;NSS\}$.  Suy ra $n(A)=3$.\\
			Xác suất của biến cố $A$ là  $\mathrm{P}(A)=\dfrac{n(A)}{n(\Omega)}=\dfrac{3}{8}$.
		\end{itemize}
	} 
\end{ex}

% nguồn: MaHoaNH_TracNghiem_Lop10
\begin{ex}%[0D0H2-2]%[ma:51a9ca8e]
	Gieo một con xúc xắc ba lần liên tiếp. Xác  suất để mặt hai chấm xuất hiện cả ba lần là
	\choice 
		{\True  $\dfrac{1}{216}$}
		 {$\dfrac{1}{72}$}
		{$\dfrac{1}{18}$} 
		 {$\dfrac{1}{20}$}
	 \loigiai{
		 Không gian mẫu là tập các bộ ba số $(i, j, k)$ với $i,j,k \in \{1,2,3,4,5,6\}$, nên
		 $
		n(\Omega) = 6 \cdot 6 \cdot  6 = 216$.\\
		Biến cố $A$: \lq\lq Cả ba lần đều xuất hiện mặt $2$ chấm\rq\rq,  chỉ có $1$ kết quả thuận lợi là $(2,2,2)$, nên $n(A) = 1$.\\
		 Vậy xác suất cần tìm là
		 $
		\mathrm{P}(A)  = \dfrac{1}{216}$.
	} 
\end{ex}

% nguồn: MaHoaNH_TracNghiem_Lop10
\begin{ex}%[0D0H2-2]%[ma:81c04b4f]
	Gieo hai  con xúc xắc. Tính xác suất của biến cố \lq\lq Tổng số chấm trên hai mặt là số lẻ\rq\rq.
	 \choice
		{$\dfrac{11}{36}$} 
		{\True $\dfrac{1}{2}$} 
		 {$\dfrac{1}{3}$}
		{$\dfrac{1}{4}$} 
	\loigiai{ 
		Không gian mẫu là các bộ $(i, j)$ với  $i,j \in \{1,2,3,4,5,6\}$.
		 \[
		n(\Omega) = 6 \cdot 6 =  36.
		\] 
		 Gọi $A$ là biến cố \lq\lq Tổng số chấm trên hai mặt là số lẻ\rq\rq.\\
		\textbf{Trường hợp 1:} Xúc xắc thứ nhất ra số lẻ ($3$ cách), xúc xắc thứ hai ra số  chẵn ($3$ cách) $\Rightarrow 3 \cdot 3 = 9$ cách.\\
		\textbf{Trường hợp 2:} Xúc xắc thứ nhất ra số chẵn ($3$ cách), xúc xắc thứ hai ra số lẻ ($3$ cách) $\Rightarrow 3 \cdot 3 =  9$ cách.\\
		$\Rightarrow n(A) = 9 + 9 = 18\Rightarrow 
		\mathrm{P}(A) = \dfrac{18}{36}  = \dfrac{1}{2}$.
	 }
\end{ex}


%%%%%%%%%%%%%%%% ĐÚNG/SAI %%%%%%%%%%%%%%%%

% nguồn: MaHoaNH_DungSai_Lop10
\begin{ex}%[0D0H2-2]%[ma:d88157e6]
	Xét phép thử $T$: \lq\lq Tung một con xúc xắc cân đối và đồng chất\rq\rq\, và biến cố $A$: \lq\lq Số  chấm xuất hiện bé hơn $5$\rq\rq.
	\choiceTF 
		{\True  Không gian mẫu $\Omega=\{1;2;3;4;5;6\}$}
		{Số kết quả thuận lợi cho  biến cố $A$ là $5$}
		{Biến cố $\overline{A}$ là biến cố \lq\lq Số chấm  xuất hiện lớn hơn $5$\rq\rq}
		 {Xác xuất của biến cố $\overline{A}$ là $\mathrm{P}(\overline{A})=\dfrac{1}{6}$}
	 \loigiai{
		\begin{itemchoice}
			\itemch  Không gian mẫu $\Omega=\{1;2;3;4;5;6\}$.
			\itemch Ta có $A=\{1;2;3;4\}$ nên có $4$ kết quả  thuận lợi cho biến cố $A$.
			\itemch Biến cố $\overline{A}$ là biến cố \lq\lq  Số chấm xuất hiện lớn hơn hoặc bằng $5$\rq\rq\,.
			 \itemch Vì $\mathrm{P}(A)=\dfrac{4}{6}=\dfrac{2}{3}$ nên $\mathrm{P}(\overline{A})=1-\dfrac{2}{3}=\dfrac{1}{3}$.
		\end{itemchoice}
	} 
\end{ex}

% nguồn: MaHoaNH_DungSai_Lop10
\begin{ex}%[0D0H2-2]%[ma:2d30d192]
	Gieo một đồng xu cân đối liên tiếp $4$ lần. Kí hiệu $S$ và $N$ tương ứng là đồng xu ra mặt sấp  và đồng xu ra mặt ngửa. Khi đó
	 \choiceTF
		{Biến cố  $A$: \lq\lq Cả $4$ lần gieo đều xuất hiện mặt ngửa\rq\rq\ là biến cố không thể}
		 {Biến cố $B$: \lq\lq Có ít nhất một lần gieo xuất hiện mặt sấp\rq\rq\ là biến cố chắc chắn}
		{\True Xác suất của biến cố $C$: \lq\lq Có đúng một lần gieo xuất hiện mặt sấp\rq\rq\  là $\mathrm{P}(C)=\dfrac{1}{4}$}
		{\True Xác suất của biến cố $D$: \lq\lq Có hai lần gieo xuất hiện mặt sấp và hai lần  gieo xuất hiện mặt ngửa\rq\rq\ là $\mathrm{P}(D)=\dfrac{3}{8}$}
	 \loigiai{
		\begin{itemchoice}
			\itemch Ta có $A=\{NNNN\}$ nên $A$ không là biến cố  không thể.
			\itemch Ta có trường hợp \lq\lq $NNNN$\rq\rq\ thuộc không gian mẫu nhưng không  thuộc $B$ nên $B$ không là biến cố chắc chắn.
			\itemch Số phần tử của không gian mẫu là $16$.\\ 
			Có $C=\{SNNN; NSNN; NNSN; NNNS\}$  nên $\mathrm{P}(C)=\dfrac{4}{16}=\dfrac{1}{4}$.
			\itemch Ta có $n(D)=\mathrm{C}_4^2=6$ nên  $\mathrm{P}(D)=\dfrac{6}{16}=\dfrac{3}{8}$.
		\end{itemchoice}
	 }
\end{ex}

% nguồn: MaHoaNH_DungSai_Lop10
\begin{ex}%[0D0H2-2]%[ma:b30706e5]
	Xét phép thử ngẫu nhiên là việc gieo hai con xúc xắc cùng một lúc.  Khi đó
	\choiceTF 
		{Số phần  tử không gian mẫu là $n(\Omega)=30$}
		 {\True Số các kết quả thuận lợi cho biến cố: \lq\lq Mặt có số chấm giống nhau xuất hiện\rq\rq\ là $6$}
		{\True Xác suất của biến cố: \lq\lq Tổng số  chấm xuất hiện trên mặt hai con xúc xắc bằng $6$\rq\rq\ là $\dfrac{5}{36}$}
		{Xác suất của biến cố: \lq\lq Tích số chấm trên hai mặt xuất hiện bằng  một số lẻ\rq\rq\ là $\dfrac{1}{5}$}
	 \loigiai{
		\begin{itemchoice}
			\itemch Số phần tử không gian mẫu là  $n(\Omega)=6\cdot 6=36$.
			\itemch Ta có $A=\{(1; 1);(2; 2);(3; 3);(4;  4);(5; 5);(6; 6)\} \Rightarrow n(A)=6$.\\
			Số các kết quả thuận  lợi cho biến cố: \lq\lq Mặt có số chấm giống nhau xuất hiện\rq\rq\ là $6$.
			\itemch Ta có $6=1+5=5+1=2+4=4+2=3+3$.\\ 
			Do đó $B=\{(1; 5);(5; 1);(2; 4);(4;  2);(3; 3)\} \Rightarrow n(B)=5$.\\
			Xác suất  biến cố $B$ là $\mathrm{P}(B)=\dfrac{n(B)}{n(\Omega)}=\dfrac{5}{36}$.
			\itemch Gọi $C$ là biến cố \lq\lq Tích  số chấm trên hai mặt xuất hiện bằng một số lẻ\rq\rq.\\
			Số các kết quả thuận lợi cho biến cố là  $n(C)=3\cdot 3=9$. Xác suất là $\mathrm{P}(C)=\dfrac{9}{36}=\dfrac{1}{4}$.
		\end{itemchoice}
	} 
\end{ex}

% nguồn: MaHoaNH_DungSai_Lop10
\begin{ex}%[0D0H2-2]%[ma:dd1b28a7]
	Tung một đồng  xu cân đối và đồng chất 3 lần.
	 \choiceTF
		 {Số phần tử của không gian mẫu là $6$}
		{\True Xác suất  để 3 lần gieo trúng mặt sấp là $\dfrac{1}{8}$ }
		{ Xác suất để hai lần nhận được mặt sấp là  $\dfrac{1}{2}$ }
		{\True Xác suất nhận  được ít nhất một mặt sấp $\dfrac{7}{8}$}
	\loigiai{ 
		\begin{itemchoice}
			 \itemch \textbf{Sai.}\\
			Ta có $\Omega=\left\{ SSS; SSN; SNS; NSS;  NNS; NSN; SNN; NNN\right\}$\\
			 Suy ra $n\left(\Omega\right)=8$.
			\itemch  \textbf{Đúng.}\\
			$A$: \lq\lq Ba lần gieo trúng  mặt sấp \rq\rq.\\
			Khi đó, $A=\left\{  SSS \right\}$ và $n( A)=1$.\\
			Suy  ra $\mathrm{P}( A)=\dfrac{n( A)}{n( \Omega)}=\dfrac{1}{8}$.
			\itemch \textbf{Sai.}\\ 
			$B$: \lq\lq Hai lần gieo trúng  mặt sấp\rq\rq.\\
			Khi đó, $A=\left\{ SSN;SNS;NSS \right\}$ và  $n( B)=3$.\\
			 Suy ra $\mathrm{P}( B)=\dfrac{n( B)}{n( \Omega)}=\dfrac{3}{8}$.
			 \itemch \textbf{Đúng.}\\
			$C$: \lq\lq Gieo được ít nhất  một mặt sấp \rq\rq\ và $\overline{C}$: \lq\lq Ba lần nhận được mặt ngửa\rq\rq.\\
			Khi đó $\mathrm{P}( C)=1-\mathrm{P}( \overline{C})=1-\dfrac{1}{8}=\dfrac{7}{8}$. 
		\end{itemchoice}
	} 
\end{ex}

% nguồn: MaHoaNH_DungSai_Lop10
\begin{ex}%[0D0H2-2]%[ma:9d674d5a]
	Gieo một đồng  tiền liên tiếp $3$ lần. Các khẳng định sau đúng hay sai?
	 \choiceTF
		{Xác suất của biến cố $3$ lần gieo xuất hiện mặt sấp là  $\dfrac{3}{8}$}
		{\True Xác suất của biến cố $3$  lần gieo như nhau là $\dfrac{1}{4}$}
		{\True Xác suất  của biến cố để ít nhất một đồng xu xuất hiện mặt ngửa là $\dfrac{7}{8}$}
		{\True Xác suất của biến cố lần đầu  xuất hiện mặt sấp là $\dfrac{1}{2}$}
	\loigiai{ 
		\begin{itemchoice}
			\itemch  Số phần tử của không gian mẫu là $n(\Omega)=2^3=8$.\\
			Gọi $ A  $ là biến cố \lq\lq $3$ lần gieo xuất hiện mặt sấp \rq\rq.\\
			Số kết quả thuận lợi cho biến cố $A$  là $\{(S,S,S)\}$ $\Rightarrow n(A)=1$.\\
			Vậy $\mathrm{P}(A)=\dfrac{n(A)}{n(\Omega)}=\dfrac{1}{8}$. 
			\itemch Số phần tử của  không gian mẫu là $n(\Omega)=2^3=8$.\\
			Gọi $ A $ là  biến cố \lq\lq $3$ lần gieo như nhau\rq\rq.\\
			Số kết quả thuận lợi cho biến cố $A$ là $\{(S,S,S);(N,N,N)\}$ $\Rightarrow n(A)=2$.\\ 
			Vậy  $\mathrm{P}(A)=\dfrac{n(A)}{n(\Omega)}=\dfrac{1}{4}$.
			\itemch  Số phần tử của không gian mẫu là $n(\Omega)=2^3=8$.\\
			Gọi $ A $ là biến cố \lq\lq ít nhất một đồng xu xuất hiện  mặt ngửa\rq\rq.\\ Suy ra $\overline{A}$ là biến cố \lq\lq Cả $3$ lần gieo đều xuất hiện mặt sấp\rq\rq.\\
			Ta có  $n\left(\overline{A}\right)=1 \Rightarrow \mathrm{P}\left(\overline{A}\right)=\dfrac{1}{8}$.
			 Vậy $\mathrm{P}(A)=1-\mathrm{P}\left(\overline{A}\right)=1-\dfrac{1}{8}=\dfrac{7}{8}$.
			\itemch Số phần tử của không gian mẫu là $n(\Omega)=2^3=8$.\\ 
			Gọi $  A $ là biến cố \lq\lq lần đầu xuất hiện mặt sấp\rq\rq.\\
			Số kết quả thuận lợi cho biến cố $A$ là $\{SSS,SSN,SNN,SNS\}$ $\Rightarrow  n(A)=4$.\\
			 Vậy $\mathrm{P}(A)=\dfrac{n(A)}{n(\Omega)}=\dfrac{1}{2}$.
		\end{itemchoice}
	 }
\end{ex}

% nguồn: MaHoaNH_DungSai_Lop10
\begin{ex}%[0D0H2-2]%[ma:0ccbb470]
	Gieo ngẫu nhiên một con súc sắc cân đối và đồng chất một lần. Các khẳng  định sau đúng hay sai?
	\choiceTF 
		{\True Xác suất của biến cố xuất hiện mặt ba chấm là  $\dfrac{1}{6}$}
		{Xác suất của biến cố xuất hiện mặt  hai chấm là $\dfrac{1}{3}$}
		{\True Xác suất của biến cố xuất hiện số  chấm lớn hơn $2$ là $\dfrac{2}{3}$}
		{Xác suất của biến cố xuất hiện  số chấm nhỏ hơn $5$ là $\dfrac{5}{6}$}
	\loigiai{ 
		\begin{itemchoice}
			\itemch Số phần tử  của không gian mẫu là $ n(\Omega)=6 $.\\
			 Gọi $ A $ là biến cố \lq\lq Xuất hiện mặt ba chấm\rq\rq.\\
			Số kết quả thuận lợi cho biến cố $  A $ là $ n(A)=1$.\\
			Vậy xác suất của biến cố $  A $ là $ \mathrm{P}(A)=\dfrac{n(A)}{n(\Omega)}=\dfrac{1}{6}$.
			\itemch Số phần tử của không gian mẫu là $  n(\Omega)=6 $.\\
			Gọi $ A $ là biến cố  \lq\lq Xuất hiện mặt hai chấm\rq\rq.\\
			Số kết quả thuận lợi cho biến cố  $ A $ là $ n(A)=1$.\\
			Vậy xác suất  của biến cố $ A $ là $ \mathrm{P}(A)=\dfrac{n(A)}{n(\Omega)}=\dfrac{1}{6}$.
			\itemch Số phần tử của không  gian mẫu là $ n(\Omega)=6 $.\\
			Gọi $ A $  là biến cố \lq\lq Xuất hiện số chấm lớn hơn $ 2 $\rq\rq.\\
			Số kết quả thuận  lợi cho biến cố $ A $ là $ n(A)=4 $.\\
			Vậy xác suất của biến cố $ A $ là $  \mathrm{P}(A)=\dfrac{n(A)}{n(\Omega)}=\dfrac{4}{6}=\dfrac{2}{3} $.
			\itemch Số phần tử của  không gian mẫu là $ n(\Omega)=6 $.\\
			Gọi $ A $ là biến  cố \lq\lq Xuất hiện số chấm nhỏ hơn $5$\rq\rq.\\
			Số kết  quả thuận lợi cho biến cố $ A $ là $ n(A)=4 $.\\
			Vậy  xác suất của biến cố $ A $ là $ \mathrm{P}(A)=\dfrac{n(A)}{n(\Omega)}=\dfrac{4}{6}=\dfrac{2}{3} $.
		\end{itemchoice}
	} 
\end{ex}

% nguồn: MaHoaNH_DungSai_Lop10
\begin{ex}%[0D0H2-2]%[ma:16c046a8]
	Gieo liên tiếp một con xúc xắc cân đối và đồng chất $3$ lần. Xét  tính Đúng/Sai của các khẳng định sau.
	\choiceTF 
		{Số  phần tử của không gian mẫu là $18$}
		{\True Xác suất để không  có lần gieo được mặt $6$ chấm là $\dfrac{125}{216}$}
		{Xác suất để có ít nhất một lần gieo được mặt $6$ chấm  là $\dfrac{75}{216}$}
		{\True Xác suất để tổng số chấm xuất hiện của $3$ lần  gieo bằng $6$ là $\dfrac{5}{108}$}
	\loigiai{ 
		Xét phép thử gieo con xúc xắc $3$ lần, số  phần tử của không gian mẫu là $n(\Omega)=6^3$.
		\begin{itemchoice}
			\itemch Số phần tử  của không gian mẫu là $6^3=216$.
			\itemch Xét biến cố $B$ \lq\lq Không có lần nào gieo được mặt $6$ chấm trong cả ba lần  gieo\rq\rq.\\
			Số khả năng thuận lợi cho biến cố $B$  là $n(B)=5^3$.\\
			 Xác suất của biến cố $B$ là $\mathrm{P}(B)=\dfrac{5^3}{6^3}=\dfrac{125}{216}$.
			\itemch Xét biến cố $C$ \lq\lq Có ít nhất một lần  gieo được mặt $6$ chấm\rq\rq.\\
			Khi đó biến cố đối của $C$ là $\overline{C}$  \lq\lq Không có lần nào gieo được mặt $6$ chấm trong cả ba lần gieo\rq\rq.\\
			Vậy số khả  năng thuận lợi cho biến cố $\overline{C}$ là $n\left(\overline{C}\right)=5^3$.\\
			Xác suất của biến cố $C$  là $\mathrm{P}(C)=1-\mathrm{P}\left(\overline{C}\right)=1-\dfrac{5^3}{6^3}=\dfrac{91}{216}$.
			\itemch Xét biến cố $D$ \lq\lq Tổng số chấm xuất  hiện của $3$ lần gieo bằng $6$ \rq\rq.
			\begin{itemize}
				\item $3$ lần gieo xuất hiện các mặt  có số chấm đều là $2$ có $1$ khả năng.
				\item $3$ lần gieo xuất hiện các mặt  có số chấm là $\{1;2;3\}$ có $3!$ khả năng.
				\item  $3$ lần gieo xuất hiện các mặt có số chấm là $\{1;1;4\}$ có $3$ khả năng.
			\end{itemize}
			Vậy số khả năng  thuận lợi cho biến cố $D$ là $n(D)=1+3!+3=10$.\\
			Xác suất của biến cố $D$  là $\mathrm{P}(D)=\dfrac{n(D)}{n(\Omega)}=\dfrac{10}{216}=\dfrac{5}{108}$.
		\end{itemchoice}
	 }
\end{ex}

% nguồn: MaHoaNH_DungSai_Lop10
\begin{ex}%[0D0H2-2]%[ma:d1b0b17c]
	Một hộp có chứa $11$ viên bi màu đỏ  và $9$ viên bi màu đen. Chọn ngẫu nhiên $4$ viên bi.
	\choiceTF 
		{\True Xác suất chọn được $4$ viên bi màu đỏ là  $\dfrac{22}{323}$}
		{\True Xác suất chọn được  số viên bi màu đỏ bằng số viên bi màu đen là $\dfrac{132}{323}$}
		 {Xác suất chọn được số viên bi màu đỏ lớn hơn số viên bi màu đen là $\dfrac{123}{323}$}
		{Xác suất chọn được ít nhất một viên bi màu đỏ là  $\dfrac{112}{323}$}
	\loigiai{ 
		 Chọn ngẫu nhiên $4$ viên bi trong hộp, ta có số phần tử của không gian mẫu là $n(\Omega)=\mathrm{C}_{20}^4=4845$.
		\begin{itemchoice}
			\itemch Gọi $A$ là biến cố \lq\lq  Chọn được $4$ viên bi màu đỏ\rq\rq.\\
			Vậy số khả năng  thuận lợi cho biến cố $A$ là $n(A)=\mathrm{C}_{11}^4=330$.\\
			Xác suất  cần tìm là $\mathrm{P}(A)=\dfrac{n(A)}{n(\Omega)}=\dfrac{330}{4845}=\dfrac{22}{323}$.
			\itemch Gọi  $B$ là biến cố \lq\lq Chọn được số viên bi màu đỏ bằng số viên bi màu đen \rq\rq.\\
			Do vậy cần chọn ra $2$ viên bi  màu đỏ và $2$ viên bi màu đen.\\
			Suy ra số khả năng thuận lợi cho biến cố $B$ là $n(B)=\mathrm{C}_{11}^2\cdot  \mathrm{C}_{9}^2=1980$.\\
			Xác suất  cần tìm là $\mathrm{P}(B)=\dfrac{n(B)}{n(\Omega)}=\dfrac{1980}{4845}=\dfrac{132}{323}$.
			 \itemch Gọi $C$ là biến cố \lq\lq Chọn được số viên bi màu đỏ lớn hơn số viên bi màu đen\rq\rq.
			\begin{itemize}
				\item Chọn được $3$ viên bi màu đỏ và $1$ viên bi màu đen có $\mathrm{C}_{11}^3\cdot  \mathrm{C}_{9}^1=1485$ cách chọn.
				\item Chọn được $4$ viên bi màu đỏ có $\mathrm{C}_{11}^4=330$  cách chọn.
			\end{itemize}
			Vậy số khả  năng thuận lợi cho biến cố $C$ là $n(C)=1485+330=1815$.\\
			Xác suất  cần tìm là $\mathrm{P}(C)=\dfrac{n(C)}{n(\Omega)}=\dfrac{1815}{4845}=\dfrac{121}{323}$.
			\itemch Gọi $D$ là biến cố \lq\lq Chọn được ít nhất một viên  bi màu đỏ\rq\rq.\\
			Khi đó biến cố đối của $D$ là $\overline{D}$ \lq\lq $4$ viên bi được chọn không có viên bi  nào màu đỏ\rq\rq.\\
			Suy ra $\overline{D}$  là biến cố \lq\lq Chọn ra $4$ viên bi đen\rq\rq.\\
			Vậy số  khả năng thuận lợi cho biến cố $\overline{D}$ là $n\left(\overline{D}\right)=\mathrm{C}_{9}^4=126$.\\
			Xác suất của biến cố  $D$ là $\mathrm{P}(D)=1-\mathrm{P}\left(\overline{D}\right)=1-\dfrac{126}{4845}=\dfrac{1573}{1615}$.
		\end{itemchoice}
	} 
\end{ex}

% nguồn: MaHoaNH_DungSai_Lop10
\begin{ex}%[0D0H2-2]%[ma:ff4b3a4f]
	Gieo một con súc sắc cân đối đồng chất hai lần. Khi đó 
	\choiceTF 
		{Số phần tử  của không gian mẫu là $12$}
		 {\True Xác suất để ít nhất một lần xuất hiện mặt sáu chấm là $\dfrac{11}{36}$}
		{Xác suất để biến cố \lq\lq  Có tổng số chấm hai mặt bằng $8$\rq\rq\ là $\dfrac{1}{6}$}
		{Xác suất để biến cố \lq\lq Có tích $2$ lần số chấm khi gieo súc  sắc là một số chẵn\rq\rq\ là $0{,}25$}
	 \loigiai{
		\begin{itemchoice}
			\itemch Số phần tử của không gian mẫu là  $n(\Omega)=6\cdot 6=36$.
			\itemch Số phần tử của không gian  mẫu là $n(\Omega)=6\cdot 6=36$.\\
			Gọi $A$ là biến cố \lq\lq Ít nhất một lần xuất hiện mặt sáu  chấm\rq\rq.\\
			Để tìm số phần tử của biến cố $A$, ta đi tìm số phần  tử của biến cố đối $\overline{A}$ là \lq\lq Không xuất hiện mặt sáu chấm\rq\rq.\\
			 Khi đó $n(\overline{A})=5\cdot 5=25\Rightarrow n(A)=36-25=11$.\\
			Vậy xác suất cần tính $\mathrm{P}(A)=\dfrac{11}{36}$. 
			 \itemch Số phần tử của không gian mẫu là $n(\Omega)=6\cdot 6=36$.\\
			Gọi $A$ là biến cố \lq\lq Số chấm trên mặt hai lần gieo có  tổng bằng $8$\rq\rq.\\
			Gọi số chấm trên mặt khi gieo lần một là $x$, số chấm trên mặt khi gieo lần hai  là $y$.\\
			Theo bài  ra, ta có \[\heva{&1\leq x\leq 6\\&1\leq y\leq 6\\&x+y=8} \Rightarrow(x; y)=\{(2; 6);(3; 5);(4; 4);(6; 2);(5; 3)\}.\]
			Khi  đó $n(A)=6$. \\
			Vậy xác suất cần  tính $\mathrm{P}(A)=\dfrac{5}{36}$.
			\itemch  Số phần tử của không gian mẫu là $n(\Omega)=6\cdot 6=36$.\\
			Gọi $A$ là biến cố  \lq\lq Tích hai lần số chấm khi gieo súc sắc là một số chẵn\rq\rq. \\
			Ta xét  các trường hợp:
			\begin{itemize}
				\item \textbf{Trường hợp 1.} Gieo lần một, số chấm xuất hiện trên mặt là số  lẻ thì khi gieo lần hai, số chấm xuất hiện phải là số chẵn. \\
				Khi đó có $3\cdot 3=9$ cách gieo. 
				\item \textbf{Trường hợp 2.}  Gieo lần một, số chấm xuất hiện trên mặt là số chẵn thì có hai trường hợp xảy ra là số chấm xuất hiện trên mặt khi gieo lần hai là số lẻ hoặc số chẵn. \\
				Khi đó  có $3\cdot 3+3\cdot 3=18$ cách gieo.
			\end{itemize}
			 Suy ra $n(A)=9+18=27$.\\
			Vậy xác suất  cần tìm là $\mathrm{P}(A)=\dfrac{27}{36}=0{,}75$.
		\end{itemchoice}
	} 
\end{ex}

% nguồn: MaHoaNH_DungSai_Lop10
\begin{ex}%[0D0V2-2]%[ma:d7d86621]
	Gieo ngẫu nhiên hai con xúc xắc  cân đối và đồng chất.
	 \choiceTF
		{\True Số phần tử của không gian mẫu $36$} 
		{\True Xác suất để sau hai lần  gieo kết quả như nhau bằng $\dfrac{1}{6}$}
		{ Xác suất để ít nhất một lần xuất  hiện mặt sáu chấm bằng $\dfrac{1}{3}$}
		{\True Xác suất để tổng số chấm  trên hai mặt bằng 7 bằng $\dfrac{1}{6}$}
	\loigiai{ 
		\begin{itemchoice}
			\itemch  \textbf{Đúng.}\\
			Số  phần tử của không gian mẫu: $n( \Omega)=6.6=36$.
			\itemch \textbf{Đúng.}\\ 
			Biến cố $A$: \lq\lq Xuất hiện hai lần gieo kết quả  như nhau \rq\rq.\\
			Ta có $A=\{(1;1);(2;2);(3;3);(4;4);(5;5);(6;6)\}$ nên $n(A)=6$.\\ 
			Suy ra $\mathrm{P}( A)=\dfrac{n( A)}{n( \Omega)}=\dfrac{6}{36}=\dfrac{1}{6}$. 
			 \itemch \textbf{Sai.}\\
			Gọi  $B$: \lq\lq Ít nhất một lần xuất hiện mặt sáu chấm\rq\rq.\\
			Khi đó $\overline{B}$: \lq\lq Không  có lần nào xuất hiện mặt sáu chấm \rq\rq.\\
			Ta có $n( {\overline{B}})=5\cdot 5=25$ nên $\mathrm{P}( B)=1-\mathrm{P} 
			(  {\overline{B}})=1-\dfrac{25}{36}=\dfrac{11}{36}$.
			 \itemch \textbf{Đúng.}\\
			 Biến cố $D$: \lq\lq Tổng số chấm trên hai mặt bằng $7$ \rq\rq.\\
			Ta có  $D=\left\{ ( 1;6);( 2;5);( 3;4);( 4;3);( 5;2);( 6;1) \right\}$ nên $n( D)=6$.\\
			Suy  ra $\mathrm{P}( D)=\dfrac{n( D)}{n( \Omega)}=\dfrac{6}{36}=\dfrac{1}{6}$.
		\end{itemchoice}
	} 
\end{ex}


%%%%%%%%%%%%%%%% TRẢ LỜI NGẮN %%%%%%%%%%%%%%%%

% nguồn: MaHoaNH_TraLoiNgan_Lop10
\begin{bt}%[0D0H2-2]%[ma:3cb1e23b]
	Xét phép thử \lq\lq Gieo một con xúc xắc (6  mặt) cân đối và đồng chất hai lần liên tiếp\rq\rq. Gọi $A$ là biến cố: \lq\lq Tổng số chấm của hai lần gieo bằng $10$\rq\rq.
	Xác suất của  biến cố $A$ bằng bao nhiêu? (Kết quả làm tròn đến hai chữ số sau dấu phẩy).
	\shortans{0,08} 
	 \loigiai{
		Không gian mẫu  $\Omega=\{(i;j) \mid i, j=1; 2; 3; 4; 5; 6\}$, trong đó $(i; j)$ là kết quả ``Lần đầu xuất hiện mặt $i$ chấm, lần sau xuất hiện mặt $j$ chấm''.\\
		Suy  ra $\mathrm{n} (\Omega)=36$.\\
		Biến cố  đề cho $A=\{(4; 6),(5; 5),(6; 4)\}$ có $\mathrm{n} (A)=3$.\\
		Vậy xác suất cần tìm là $\mathrm{P}(A)=\dfrac{\mathrm{n} (A)}{\mathrm{n} (\Omega)}=\dfrac{1}{12} \approx 0{,}08$. 
	 }
\end{bt}

% nguồn: MaHoaNH_TraLoiNgan_Lop10
\begin{bt}%[0D0V2-2]%[ma:25c1471b]
	Gieo đồng thời hai viên xúc xắc cân đối và đồng chất. Tính xác suất để tổng số chấm xuất hiện trên hai viên xúc xắc bằng $9$. (kết quả làm tròn đến hàng phần  trăm)
	\shortans{0,11} 
	\loigiai{ 
		Số phần  tử của không gian mẫu là $n(\Omega)=36$.\\
		Gọi $A$ là biến cố  tổng số chấm trên hai viên xúc xắc bằng $9$.\\
		$A=\{(3; 6),(4; 5);(5; 4);(6; 3)\}$. Do đó,  ta có $n(A)=4$.\\
		Vậy  xác suất của biến cố $A$ là $\mathrm{P}(A)=\dfrac{n(A)}{n(\Omega)}=\dfrac{4}{36}=\dfrac{1}{9}\approx0{,}11$.
	 }
\end{bt}

% nguồn: MaHoaNH_TraLoiNgan_Lop10
\begin{bt}%[0D0V2-2]%[ma:9ebcaf3e]
	Gieo một viên xúc  xắc $6$ mặt cân đối và đồng chất liên tiếp năm lần. Tính xác suất để mặt $6$ chấm xuất hiện ít nhất một lần (kết quả làm tròn đến hàng phần mười).
	 \shortans{0,6}
	\loigiai{ 
		Gọi $A$ là  biến cố \lq\lq Mặt $6$ chấm không xuất hiện lần nào\rq\rq.\\
		Suy ra $\overline{A}$ là biến cố \lq\lq  Mặt $6$ chấm xuất hiện ít nhất một lần\rq\rq.\\
		 Ta có $n(\Omega)=6^5=7776$, $n(A)=5^3=3125$.\\
		Do đó, xác suất của biến cố $A$  là $\mathrm{P}(A)=\dfrac{n(A)}{n(\Omega)}=\dfrac{3125}{7776}$.\\
	Vậy xác suất của biến cố \lq\lq Mặt $6$ chấm xuất  hiện ít nhất một lần\rq\rq \, là $\mathrm{P}(\overline{A})=1-\dfrac{3125}{7776}=\dfrac{4651}{7776}\approx 0{,}6$.}
\end{bt}

% nguồn: MaHoaNH_TraLoiNgan_Lop10
\begin{bt}%[0D0V2-2]%[ma:d19b7307]
	Gieo một viên xúc xắc 6 mặt cân đối và đồng  chất liên tiếp năm lần. Tính xác suất để mặt 6 chấm xuất hiện ít nhất một lần.
	 \shortans{0,6}
	 \loigiai{
		Gọi $A$ là biến cố \lq\lq Mặt 6 chấm không xuất hiện lần nào\rq\rq. Suy ra $\overline{A}$ là biến cố \lq\lq Mặt 6 chấm xuất hiện ít  nhất một lần\rq\rq.\\
		Số phần tử của không gian mẫu $n(\Omega)=6 \cdot 6 \cdot 6 \cdot 6  \cdot 6=7\,776$.
		\\Số phần tử của  biến cố $A$ là $n(A)=5 \cdot 5 \cdot 5 \cdot 5 \cdot 5=3\,125$.\\
		Do  đó, xác suất của biến cố $A$ là $\mathrm{P}(A)=\dfrac{n(A)}{n(\Omega)}=\dfrac{3\,125}{7\,776}$.\\
		Vậy xác suất  của biến cố $A$ là $\mathrm{P}\left(\overline{A}\right)=1-\dfrac{3\,125}{7\,776}=\dfrac{4\,651}{7\,776}$.
	} 
\end{bt}


%%%%%%%%%%%%%%%% TỰ LUẬN %%%%%%%%%%%%%%%%

% nguồn: MaHoaNH_TuLuan_Lop10
\begin{bt}%[0D0N2-2]%[ma:29137167]
	Gieo lần lượt hai con  xúc xắc cân đối đồng chất. Tính xác suất của biến cố \lq \lq Số chấm xuất hiện trên mặt hai con xúc xắc giống nhau \rq \rq.
	\loigiai{Số phần  tử của không gian mẫu là $n(\Omega)=6 \cdot 6=36$.\\
		Gọi $A$ là biến cố: \lq \lq Số chấm xuất  hiện trên mặt hai con xúc xắc giống nhau\rq \rq.\\
		Có $6$ kết quả thuận  lợi cho biến cố $A$ là $\left\{1;1\right\}$; $\left\{2;2\right\}$; $\left\{3;3\right\}$; $\left\{4;4\right\}$; $\left\{5;5\right\}$; $\left\{6;6\right\}$.\\
		Xác suất để số chấm xuất hiện trên mặt hai con xúc xắc giống nhau là $$\mathrm{P}(A)=\dfrac{n(A)}{n(\Omega)}=\dfrac{6}{36}=\dfrac{1}{6}.$$ 
	 }
\end{bt}

% nguồn: MaHoaNH_TuLuan_Lop10
\begin{bt}%[0D0N2-2]%[ma:ddd1826a]
	Gieo một con xúc xắc cân đối đồng chất. Tính xác suất của biến  cố \lq \lq Mặt có số chấm chẵn xuất hiện\rq \rq.
	\loigiai{Không gian mẫu $\Omega = \left\{1;2;3;4;5;6\right\}$  suy ra $n(\Omega)=6$.\\
		Gọi $A$ là biến cố: \lq \lq Mặt có số chấm chẵn xuất hiện\rq \rq. 
		 \\
		Khi đó $A=\left\{2;4;6\right\}$ nên  $n(A)=3$.\\
		Do đó $\mathrm{P}(A)=\dfrac{n(A)}{n(\Omega)}=\dfrac{1}{2}$. 
	 }
\end{bt}

% nguồn: MaHoaNH_TuLuan_Lop10
\begin{bt}%[0D0H2-2]%[ma:d97e261e]
	Gieo một con xúc xắc cân đối đồng chất $2$ lần, tính xác suất của biến cố \lq Tích số  chấm ở $2$ lần khi gieo xúc xắc là một số chẵn\rq \rq.
	\loigiai{Số phần tử của  không gian mẫu là $n(\Omega)=6 \cdot 6=36$.\\
		Gọi $A$ là biến cố: \lq \lq Tích hai lần số chấm khi gieo  xúc xắc là một số chẵn \rq \rq.\\
		Ta  xét các trường hợp sau:\\
		\textbf{Trường hợp 1:} Gieo lần một, số chấm xuất hiện  trên mặt là số lẻ thì khi gieo lần hai, số chấm xuất hiện phải là số chẵn. Khi đó có $3 \cdot 3=9$ cách gieo.\\
		\textbf{Trường hợp 2}: Gieo lần một, số chấm xuất hiện trên mặt  là số chẵn thì có hai trường hợp xảy ra là số chấm xuất hiện trên mặt khi gieo lần hai là số lẻ hoặc số chẵn. \\
		 Khi đó có $3 \cdot 3+3 \cdot 3 =18$ cách gieo.\\
		Suy ra số kết quả thuận lợi cho biến cố là  $n(A)=9+18=27$.\\
		Vậy xác suất cần tìm là  $\mathrm{P}(A)=\dfrac{27}{36}=0{,}75$.
	} 
\end{bt}

% nguồn: MaHoaNH_TuLuan_Lop10
\begin{bt}%[0D0K2-2]%[ma:c86eb43c]
	Chọn ngẫu nhiên $ 2 $ số  nguyên dương không vượt quá $ 20 $. Tính xác suất để chọn được $ 2 $ số có tích là một số chẵn.
	\loigiai{ 
		Số các kết quả có thể xảy ra là  $ n(\Omega)= \mathrm{C}_{20}^2=190$.\\
		Gọi $ C $ là biến cố “2 số có tích là  một số chẵn” nên
		$  n(C)=\mathrm{C}_{10}^2+\mathrm{C}_{10}^1\cdot \mathrm{C}_{10}^1=145 $ cách.\\
		Do đó $ \mathrm{P}(C)=\dfrac{\mathrm{C}_{10}^2+\mathrm{C}_{10}^1\cdot  \mathrm{C}_{10}^1}{\mathrm{C}_{20}^2}=\dfrac{29}{38} $.
	 }
\end{bt}

% nguồn: MaHoaNH_TuLuan_Lop10
\begin{bt}%[0D0K2-2]%[ma:2cdbcedd]
	Trong hộp có  $ 3 $ viên bi xanh và $ 5 $ viên bi đỏ. Lấy ngẫu nhiên trong hộp $ 3 $ viên bi. Tính xác suất của biến cố: \lq\lq Lấy ra được $ 3 $ viên bi màu đỏ \rq\rq.
	\loigiai{ 
	Ta có $ n(\Omega)  =\mathrm{C}_8^3=56$.\\
	Gọi A là biến cố ''Lấy ra được $  3 $ viên bi màu đỏ''.\\
	Khi  đó $ n(A)=\mathrm{C}_5^3=10 $.\\
	Vậy  $ P(A)=\dfrac{n(A)}{n(\Omega)}=\dfrac{10}{56}=\dfrac{5}{28} $.
	} 
\end{bt}

% nguồn: MaHoaNH_TuLuan_Lop10
\begin{bt}%[0D0B2-2]%[ma:f8c04d32]
	Một hộp có $ 3 $ bi trắng, $ 4 $ bi xanh, $ 5 $ bi vàng. Lấy ngẫu nhiên từ hộp ra $ 3 $ viên bi. Tính xác suất  của mỗi biến cố sau
	\begin{enumerate}
		\item $ A $: “Ba viên bi  lấy ra có cùng màu”.
		\item $ B $: “Ba viên bi lấy ra không có viên bi màu trắng”. 
	\end{enumerate}
	\loigiai{ 
		\begin{enumerate}
			\item $\mathrm{P}(A)=\dfrac{\mathrm{C}_3^3+\mathrm{C}_4^3+\mathrm{C}_5^3}{\mathrm{C}_{12}^3}=\dfrac{3}{44}  $.
			\item $ \mathrm{P}(B)=\dfrac{\mathrm{C}_9^3}{\mathrm{C}_{12}^3}=\dfrac{7}{22}  $.
		\end{enumerate}
	 }
\end{bt}

% nguồn: MaHoaNH_TuLuan_Lop10
\begin{bt}%[0D0B2-2]%[ma:4338c65a]
	Chọn ngẫu nhiên một  số tự nhiên có $5$ chữ số khác nhau. Tính xác suất để chọn được số chia hết cho $25$.
	\loigiai{ 
		\begin{itemize}
			\item Số các số tự nhiên có $5$  chữ số khác nhau là $9\cdot \mathrm{A}_{9}^{4}=27216$.
			\item  Số tự nhiên có $5$ chữ số khác nhau và chia hết cho $25$ có dạng $\overline{abc25}$, $\overline{abc50}$, $\overline{abc75}$.\\
			Có $7\cdot \mathrm{A}_{7}^{2}\cdot 2$  số dạng $\overline{abc25}$, $\overline{abc75}$.\\
			Có $\mathrm{A}_{8}^{3}$  số dạng $\overline{abc50}$.\\
			Suy ra số các số tự nhiên có $5$ chữ số khác  nhau và chia hết cho $25$ là $7\cdot \mathrm{A}_{7}^{2}\cdot 2+\mathrm{A}_{8}^{3}=924$.
			\item Phép thử $T\colon$ \lq\lq Chọn  ngẫu nhiên một số tự nhiên có $5$ chữ số khác nhau\rq\rq.\\
			Có $n(\Omega)=\mathrm{C}_{27216}^{1}=27216$. 
			\item Xét biến cố $A\colon$ \lq\lq Chọn được số chia hết cho  $25$\rq\rq.\\
			Có  $n(A)=\mathrm{C}_{924}^{1}=924$.
			\item Xác suất cần tìm  là $\mathrm{P}(A)=\dfrac{n(A)}{n(\Omega)}=\dfrac{924}{27216}=\dfrac{11}{324}$.
		\end{itemize}
	} 
\end{bt}
